% !TEX program = xelatex
\documentclass[12pt]{article}
\usepackage[UTF8]{ctex}
\usepackage{booktabs}
\usepackage{amsmath}
\usepackage{amssymb}
\usepackage{geometry}
\usepackage{caption}
\usepackage{threeparttable}
\usepackage{float}
\usepackage{setspace}
\usepackage{pdflscape}
\usepackage{array}
\usepackage{enumitem}
\usepackage{hyperref}

\geometry{left=1.1in, right=1.1in, top=1in, bottom=1in}
\setstretch{1.2}
\captionsetup{labelfont=bf, font=small}

\newcommand{\sym}[1]{\ensuremath{^{\text{#1}}}}
\newcommand{\se}[1]{{\small(#1)}}
\newcommand{\abs}{\textemdash}
\newcommand{\mis}{}
\newcolumntype{L}[1]{>{\raggedright\arraybackslash}p{#1}}

% -------------------------------------------------------
\begin{document}

\title{\textbf{外包采购价格决定因素：回归结果与数据说明}}
\author{}
\date{\today}
\maketitle



% ============================================================
\section{回归结果}
% ============================================================

\subsection*{样本概览}

本报告从VAT数据库中随机抽取3,410家企业。样本经历以下筛选步骤：

\begin{itemize}[nosep]
  \item 3,410家企业中，3,376家在采购端有9位产品码记录，其他企业的产出为乱码；
  \item 外包定义要求同一企业对同一产品同时有采购记录和销售记录，满足条件的企业为2,108家；
  \item 进一步剔除233家中介企业（外包比例$>90\%$），最终回归样本包含1,875家企业。
\end{itemize}

被解释变量为外包采购单价对数：$\ln p_{fjct}$，其中 $f$ 为企业，$j$ 为9位商品码，$c$ 为城市，$t$ 为年份（2017年）。

回归中涉及三类解释变量：
\begin{itemize}[nosep]
  \item 企业层面（$x$）：$\ln\mathrm{Output}_{ft}$（企业年总增值税销售额对数，规模代理）；$\ln K_{ft}$（企业资产总额对数）；$\ln Q_{fjt}$（企业对该产品的销售/产出数量对数，代理产量）。前两者为 firm$\times$year 级，单年截面下被 Firm FE 吸收；$\ln Q_{fjt}$ 为 firm$\times$product 级，FE 规格下仍可识别。
  \item 市场层面（$z$）：$\ln n^B_{jct}$（城市-产品买方数对数，全量口径）；$\ln n^S_{jct}$（卖方数对数）；$\ln q^m_{jct}$（市场总采购量对数）；$\ln\bar{p}^{B}_{jct}$（买方侧 leave-one-out 市场均价对数）；$\ln\bar{p}^{S}_{jct}$（卖方侧 leave-one-out 市场均价对数）。两个 LOO 均价均从城市-产品总量中剔除本企业自身的买方/卖方交易，以消除被解释变量与控制变量之间的机械相关。
  \item 产品相似性（$s$）：$S_{mj}$（投入结构余弦相似度）；$C_{mj}$（产出结构相似度）。
\end{itemize}

从产品结构来看，采购端3,376家企业平均采购约125种产品。回归面板中1,875家企业平均外包产品数约25种（均值25.0，中位数3）。

样本规模说明：含资本变量（$\ln K_{ft}$）的规格有效观测数约29,000（覆盖率约62.3\%）；不含资本变量的规格有效观测数约45,700–46,400。所有回归均在企业层面聚类稳健标准误，固定效应规格包含企业（$\delta_f$）、产品（$\delta_j$）、城市（$\delta_c$）固定效应。

\subsection*{识别说明}

由于数据为单年截面（2017年），企业层面变量 $\ln\mathrm{Output}_{ft}$、$\ln K_{ft}$ 与 Firm FE 完全共线，仅在 OLS 规格可识别。产量 $\ln Q_{fjt}$ 在 firm$\times$product 层面变化，即使加入 Firm FE 仍可识别——数量折扣效应贯穿所有规格稳健显著为负。相似度 $S_{mj}$、$C_{mj}$ 同样在 firm$\times$product 层面，Firm FE 下可识别，但 Product FE 下大幅衰减。市场变量在 product$\times$city 层面变化，Firm FE 下可识别。

% ------------------------------------------------------------
\subsection{仅企业层面变量}
% ------------------------------------------------------------

\subsubsection*{回归方程}

\begin{equation}
  \ln p_{fjct} = \gamma_1 \ln\mathrm{Output}_{ft} + \gamma_2 \ln K_{ft}
    + \gamma_3 \ln Q_{fjt}
    + \delta_f + \delta_j + \delta_c + \varepsilon_{fjct}
\end{equation}

\subsubsection*{变量说明}

\begin{table}[H]
  \centering\small
  \begin{tabular}{L{3.5cm} L{10cm}}
    \toprule
    变量 & 说明 \\
    \midrule
    $p_{fjct}$ & 因变量，外包采购单价对数 \\[3pt]
    $\ln\mathrm{Output}_{ft}$ & 企业年销售额对数，规模代理 \\[3pt]
    $\ln K_{ft}$ & 企业资产总额对数，覆盖约62.3\%样本企业 \\[3pt]
    $\ln Q_{fjt}$ & 企业对该产品的销售（产出）数量对数，代理产量 \\[3pt]
    $\delta_f,\ \delta_j,\ \delta_c$ & 企业、产品、城市固定效应 \\
    \bottomrule
  \end{tabular}
\end{table}

\subsubsection*{关键发现}

\begin{itemize}[nosep]
  \item OLS 中，企业规模 $\ln\mathrm{Output}$ 显著为正（$+0.234^{***}$）：规模越大、外包采购单价越高；资本 $\ln K$ 显著为负（$-0.156^{***}$）：资本越雄厚、采购单价越低。
  \item 产量 $\ln Q$ 在 OLS 与加入 Firm FE 后均显著为负（$-0.177^{***}$、$-0.169^{***}$），即产量越大、外包采购单价越低，体现数量折扣。
\end{itemize}

\subsubsection*{回归结果}

\begin{table}[H]
  \centering
  \caption{T1：仅企业层面变量}
  \begin{threeparttable}
    \begin{tabular}{lcc}
      \toprule
       & (1) OLS & (2) +Firm FE \\
      \midrule
      $\ln\mathrm{Output}$ & $0.234\sym{***}$ & \abs \\
                           & \se{0.051}       &     \\[4pt]
      $\ln K$              & $-0.156\sym{***}$ & \abs \\
                           & \se{0.047}        &     \\[4pt]
      $\ln Q$              & $-0.177\sym{***}$ & $-0.169\sym{***}$ \\
                           & \se{0.013}        & \se{0.013} \\[4pt]
      Cons               & $3.527\sym{***}$ & $4.969\sym{***}$ \\
                           & \se{0.452}       & \se{0.081}       \\
      \midrule
      Firm FE              & No  & Yes \\
      Product FE           & No  & No  \\
      City FE              & No  & No  \\
      $N$                  & 29,247 & 28,915 \\
      Adj $R^2$            & 0.081  & 0.311  \\
      \bottomrule
    \end{tabular}

  \end{threeparttable}
\end{table}

% ------------------------------------------------------------
\subsection{企业层面 + 产品相似性}
% ------------------------------------------------------------

\subsubsection*{回归方程}

\begin{equation}
  \ln p_{fjct} = \gamma_1 \ln\mathrm{Output}_{ft} + \gamma_2 \ln K_{ft}
    + \gamma_3 \ln Q_{fjt}
    + \beta_1 S_{mj} + \beta_2 C_{mj}
    + \delta_f + \delta_j + \delta_c + \varepsilon_{fjct}
\end{equation}

\subsubsection*{变量说明}

\begin{table}[H]
  \centering\small
  \begin{tabular}{L{3.5cm} L{10cm}}
    \toprule
    变量 & 说明 \\
    \midrule
    $S_{mj}$ & 投入结构相似度\\[3pt]
    $C_{mj}$ & 产出结构相似度 \\[3pt]
    \bottomrule
  \end{tabular}
\end{table}

\subsubsection*{关键发现}

\begin{itemize}[nosep]
  \item OLS 中企业规模显著为正（$+0.267^{***}$）、资本显著为负（$-0.112^{***}$），与 T1 一致。
  \item 相似度 $S_{mj}$、$C_{mj}$ 在 OLS 强显著为正（$+1.491^{***}$、$+1.481^{***}$），加入 Product FE 后大幅衰减至勉强显著（$0.186^{*}$、$0.263^{*}$），即相似度效应主要由产品层异质性解释。
  \item 产量 $\ln Q$ 从 OLS（$-0.226^{***}$）到全 FE（$-0.070^{***}$）始终显著为负。
\end{itemize}

\subsubsection*{回归结果}

\begin{table}[H]
  \centering
  \caption{T2：企业层面 + 产品相似性}
  \begin{threeparttable}
    \begin{tabular}{lcccc}
      \toprule
       & (1) OLS & (2) +Firm FE & (3) +Firm+Prod & (4) +Firm+Prod+City \\
      \midrule
      $\ln\mathrm{Output}$ & $0.267\sym{***}$ & \abs & \abs & \abs \\
                           & \se{0.042}       &     &     &     \\[4pt]
      $\ln K$              & $-0.112\sym{***}$ & \abs & \abs & \abs \\
                           & \se{0.040}        &     &     &     \\[4pt]
      $\ln Q$              & $-0.226\sym{***}$ & $-0.193\sym{***}$ & $-0.070\sym{***}$ & $-0.070\sym{***}$ \\
                           & \se{0.013}        & \se{0.014}        & \se{0.008}        & \se{0.008} \\[4pt]
      $S_{mj}$             & $1.491\sym{***}$ & $0.460\sym{**}$ & $0.186\sym{*}$ & $0.186\sym{*}$ \\
                           & \se{0.255}       & \se{0.186}      & \se{0.109}     & \se{0.109} \\[4pt]
      $C_{mj}$             & $1.481\sym{***}$ & $1.361\sym{***}$ & $0.263\sym{*}$ & $0.263\sym{*}$ \\
                           & \se{0.285}       & \se{0.210}       & \se{0.141}     & \se{0.141} \\[4pt]
      Cons               & $1.826\sym{***}$ & $4.846\sym{***}$ & $4.274\sym{***}$ & $4.274\sym{***}$ \\
                           & \se{0.452}       & \se{0.076}       & \se{0.047}       & \se{0.048}       \\
      \midrule
      Firm FE              & No  & Yes & Yes & Yes \\
      Product FE           & No  & No  & Yes & Yes \\
      City FE              & No  & No  & No  & Yes \\
      $N$                  & 29,237 & 28,905 & 28,552 & 28,552 \\
      Adj $R^2$            & 0.127  & 0.321  & 0.595  & 0.592  \\
      \bottomrule
    \end{tabular}
  \end{threeparttable}
\end{table}

% ------------------------------------------------------------
\subsection{仅市场层面变量}
% ------------------------------------------------------------

\subsubsection*{回归方程}

\begin{equation}
  \ln p_{fjct} = \lambda_1 \ln n^B_{jct} + \lambda_2 \ln n^S_{jct}
    + \lambda_3 \ln q^m_{jct} + \lambda_4 \ln\bar{p}^{B}_{jct} + \lambda_5 \ln\bar{p}^{S}_{jct}
    + \delta_f + \delta_j + \delta_c + \varepsilon_{fjct}
\end{equation}

\subsubsection*{变量说明}

\begin{table}[H]
  \centering\small
  \begin{tabular}{L{3.5cm} L{10cm}}
    \toprule
    变量 & 说明 \\
    \midrule
    $\ln n^B_{jct}$       & 城市 $c$ 中产品 $j$ 的买方数对数，需求侧市场厚度 \\[3pt]
    $\ln n^S_{jct}$       & 卖方数对数，供给侧竞争 \\[3pt]
    $\ln q^m_{jct}$       & 市场总采购量对数 \\[3pt]
    $\ln\bar{p}^{B}_{jct}$  & 买方侧 leave-one-out 市场均价对数（已剔除本企业自身） \\[3pt]
    $\ln\bar{p}^{S}_{jct}$  & 卖方侧 leave-one-out 市场均价对数（已剔除本企业自身） \\
    \bottomrule
  \end{tabular}
\end{table}

\subsubsection*{关键发现}

\begin{itemize}[nosep]
  \item 两个 leave-one-out 均价的对比是核心发现：买方侧 $\ln\bar{p}^{B}$ 在 OLS/Firm FE 强显著为正（$0.440^{***}$、$0.365^{***}$），但加入 Product FE 后塌为 $-0.029^{*}$，说明此前强正相关主要来自产品构成与机械相关。
  \item 而卖方侧 $\ln\bar{p}^{S}$ 即使在全 FE 下仍显著为正（$0.052^{***}$），携带了 Product FE 之外的独立价格信息。
  \item 买方数 $\ln n^B$ 全规格稳健为正（$0.115^{***}$）；市场采购量 $\ln q^m$ 在 Product FE 下显著为负（$-0.100^{***}$）。
\end{itemize}

\subsubsection*{回归结果}

\begin{table}[H]
  \centering
  \caption{T3：仅市场层面变量}
  \begin{threeparttable}
    \begin{tabular}{lcccc}
      \toprule
       & (1) OLS & (2) +Firm FE & (3) +Firm+Prod & (4) +Firm+Prod+City \\
      \midrule
      $\ln n^B$             & $0.170\sym{***}$ & $0.098\sym{***}$ & $0.115\sym{***}$ & $0.115\sym{***}$ \\
                            & \se{0.051}       & \se{0.029}       & \se{0.040}       & \se{0.040} \\[4pt]
      $\ln n^S$             & $0.104\sym{*}$  & $0.064\sym{*}$  & $-0.006$ & $-0.006$ \\
                            & \se{0.061}      & \se{0.036}      & \se{0.035} & \se{0.035} \\[4pt]
      $\ln q^m$             & $-0.009$ & $-0.029\sym{*}$ & $-0.100\sym{***}$ & $-0.100\sym{***}$ \\
                            & \se{0.017} & \se{0.015}      & \se{0.017}        & \se{0.017} \\[4pt]
      $\ln\bar{p}^{B}$      & $0.440\sym{***}$ & $0.365\sym{***}$ & $-0.029\sym{*}$ & $-0.029\sym{*}$ \\
                            & \se{0.017}       & \se{0.016}       & \se{0.017} & \se{0.017} \\[4pt]
      $\ln\bar{p}^{S}$      & $0.302\sym{***}$ & $0.246\sym{***}$ & $0.052\sym{***}$ & $0.052\sym{***}$ \\
                            & \se{0.020}       & \se{0.016}       & \se{0.011} & \se{0.011} \\[4pt]
      Cons                & $0.218$ & $1.643\sym{***}$ & $4.912\sym{***}$ & $4.912\sym{***}$ \\
                            & \se{0.259}     & \se{0.208}       & \se{0.302}       & \se{0.302}       \\
      \midrule
      Firm FE               & No  & Yes & Yes & Yes \\
      Product FE            & No  & No  & Yes & Yes \\
      City FE               & No  & No  & No  & Yes \\
      $N$                   & 46,444 & 45,956 & 45,669 & 45,669 \\
      Adj $R^2$             & 0.310  & 0.468  & 0.604  & 0.602  \\
      \bottomrule
    \end{tabular}
  \end{threeparttable}
\end{table}

% ------------------------------------------------------------
\subsection{企业层面 + 市场层面}
% ------------------------------------------------------------

\subsubsection*{回归方程}

\begin{equation}
  \ln p_{fjct} = \boldsymbol{\gamma} \cdot \mathbf{x}_{ft}
    + \boldsymbol{\lambda} \cdot \mathbf{z}_{jct}
    + \delta_f + \delta_j + \delta_c + \varepsilon_{fjct}
\end{equation}

其中 $\mathbf{x}$ 为企业层面变量（含产量 $\ln Q_{fjt}$），$\mathbf{z}_{jct}$ 为市场层面变量。

\subsubsection*{关键发现}

\begin{itemize}[nosep]
  \item OLS 中企业规模为正（$+0.157^{***}$）、资本为负（$-0.101^{***}$）；产量 $\ln Q$ 从 OLS（$-0.126^{***}$）到全 FE（$-0.061^{***}$）始终显著为负。
  \item 买方侧 $\ln\bar{p}^{B}$ 在 OLS/Firm FE 显著为正（$0.503^{***}$、$0.403^{***}$），加入 Product FE 后归零（$-0.008$）；卖方侧 $\ln\bar{p}^{S}$ 在全 FE 下仍显著为正（$0.047^{***}$）。
  \item 市场采购量 $\ln q^m$ 在 OLS 为正（$0.077^{***}$）、全 FE 下转负（$-0.078^{***}$）；买方数在全 FE 下为正（$0.106^{**}$）。
\end{itemize}



\subsubsection*{回归结果}

\begin{table}[H]
  \centering
  \caption{T4：企业层面 + 市场层面}
  \begin{threeparttable}
    \begin{tabular}{lcccc}
      \toprule
       & (1) OLS & (2) +Firm FE & (3) +Firm+Prod & (4) +Firm+Prod+City \\
      \midrule
      $\ln\mathrm{Output}$ & $0.157\sym{***}$ & \abs & \abs & \abs \\
                           & \se{0.037}       &     &     &     \\[4pt]
      $\ln K$              & $-0.101\sym{***}$ & \abs & \abs & \abs \\
                           & \se{0.033}        &     &     &     \\[4pt]
      $\ln Q$              & $-0.126\sym{***}$ & $-0.122\sym{***}$ & $-0.061\sym{***}$ & $-0.061\sym{***}$ \\
                           & \se{0.010}        & \se{0.009}        & \se{0.008}        & \se{0.008} \\[4pt]
      $\ln n^B$            & $0.045$ & $0.132\sym{***}$ & $0.106\sym{**}$  & $0.106\sym{**}$  \\
                           & \se{0.060}      & \se{0.032}       & \se{0.052} & \se{0.052} \\[4pt]
      $\ln n^S$            & $0.138\sym{*}$  & $-0.008$  & $0.012$  & $0.012$  \\
                           & \se{0.075} & \se{0.040} & \se{0.040} & \se{0.040} \\[4pt]
      $\ln q^m$            & $0.077\sym{***}$ & $0.033\sym{**}$ & $-0.078\sym{***}$ & $-0.078\sym{***}$ \\
                           & \se{0.021} & \se{0.016}       & \se{0.022}        & \se{0.022} \\[4pt]
      $\ln\bar{p}^{B}$     & $0.503\sym{***}$ & $0.403\sym{***}$ & $-0.008$  & $-0.008$  \\
                           & \se{0.023}       & \se{0.019}       & \se{0.021} & \se{0.021} \\[4pt]
      $\ln\bar{p}^{S}$     & $0.242\sym{***}$ & $0.219\sym{***}$ & $0.047\sym{***}$  & $0.047\sym{***}$  \\
                           & \se{0.022}       & \se{0.019}       & \se{0.015} & \se{0.015} \\[4pt]
      Cons               & $-0.938\sym{*}$  & $1.385\sym{***}$ & $4.664\sym{***}$ & $4.664\sym{***}$ \\
                           & \se{0.503} & \se{0.241}      & \se{0.363}       & \se{0.364}       \\
      \midrule
      Firm FE              & No  & Yes & Yes & Yes \\
      Product FE           & No  & No  & Yes & Yes \\
      City FE              & No  & No  & No  & Yes \\
      $N$                  & 28,967 & 28,636 & 28,287 & 28,287 \\
      Adj $R^2$            & 0.332  & 0.472  & 0.597  & 0.595  \\
      \bottomrule
    \end{tabular}
  \end{threeparttable}
\end{table}



% ------------------------------------------------------------
\subsection{全规格}
% ------------------------------------------------------------

\subsubsection*{回归方程}

\begin{equation}
  \ln p_{fjct} = \boldsymbol{\gamma} \cdot \mathbf{x}_{ft}
    + \boldsymbol{\beta} \cdot \mathbf{s}_{fj}
    + \boldsymbol{\lambda} \cdot \mathbf{z}_{jct}
    + \delta_f + \delta_j + \delta_c + \varepsilon_{fjct}
\end{equation}

\subsubsection*{关键发现}

\begin{itemize}[nosep]
  \item OLS 中企业规模为正（$+0.181^{***}$）、资本为负（$-0.074^{**}$）；相似度 $S_{mj}$、$C_{mj}$ 在 OLS 强显著为正（$1.076^{***}$、$0.888^{***}$），Product FE 下衰减至勉强显著（$0.197^{*}$、$0.258^{*}$）。
  \item 买方侧 $\ln\bar{p}^{B}$ 在全 FE 下归零（$-0.011$），但卖方侧 $\ln\bar{p}^{S}$ 仍显著为正（$0.047^{***}$）。
  \item 买方数 $\ln n^B$（$0.102^{**}$）、市场采购量 $\ln q^m$（$-0.082^{***}$）与产量 $\ln Q$（$-0.068^{***}$）在全 FE 下均保留信息。
\end{itemize}



\subsubsection*{回归结果}

\begin{table}[H]
  \centering
  \caption{T5：全规格（企业层面 + 产品相似性 + 市场层面）}
  \begin{threeparttable}
    \begin{tabular}{lcccc}
      \toprule
       & (1) OLS & (2) +Firm FE & (3) +Firm+Prod & (4) +Firm+Prod+City \\
      \midrule
      $\ln\mathrm{Output}$ & $0.181\sym{***}$ & \abs & \abs & \abs \\
                           & \se{0.033}       &     &     &     \\[4pt]
      $\ln K$              & $-0.074\sym{**}$ & \abs & \abs & \abs \\
                           & \se{0.029}        &     &     &     \\[4pt]
      $\ln Q$              & $-0.158\sym{***}$ & $-0.138\sym{***}$ & $-0.068\sym{***}$ & $-0.068\sym{***}$ \\
                           & \se{0.010}        & \se{0.009}        & \se{0.008}        & \se{0.008} \\[4pt]
      $S_{mj}$             & $1.076\sym{***}$ & $0.386\sym{***}$  & $0.197\sym{*}$  & $0.197\sym{*}$  \\
                           & \se{0.185}       & \se{0.115} & \se{0.107} & \se{0.108} \\[4pt]
      $C_{mj}$             & $0.888\sym{***}$  & $0.795\sym{***}$  & $0.258\sym{*}$ & $0.258\sym{*}$ \\
                           & \se{0.229} & \se{0.149} & \se{0.139} & \se{0.140} \\[4pt]
      $\ln n^B$            & $0.092$ & $0.135\sym{***}$ & $0.102\sym{**}$  & $0.102\sym{**}$  \\
                           & \se{0.059}       & \se{0.031}       & \se{0.052} & \se{0.052} \\[4pt]
      $\ln n^S$            & $0.107$  & $-0.003$  & $0.015$  & $0.015$  \\
                           & \se{0.074} & \se{0.041} & \se{0.040} & \se{0.040} \\[4pt]
      $\ln q^m$            & $0.048\sym{**}$ & $0.025$ & $-0.082\sym{***}$ & $-0.082\sym{***}$ \\
                           & \se{0.020}      & \se{0.016}       & \se{0.022}        & \se{0.022} \\[4pt]
      $\ln\bar{p}^{B}$     & $0.467\sym{***}$ & $0.393\sym{***}$ & $-0.011$  & $-0.011$  \\
                           & \se{0.022}       & \se{0.019}       & \se{0.021} & \se{0.021} \\[4pt]
      $\ln\bar{p}^{S}$     & $0.228\sym{***}$ & $0.215\sym{***}$ & $0.047\sym{***}$  & $0.047\sym{***}$  \\
                           & \se{0.021}       & \se{0.019}       & \se{0.015} & \se{0.015} \\[4pt]
      Cons               & $-1.602\sym{***}$ & $1.413\sym{***}$ & $4.712\sym{***}$ & $4.712\sym{***}$ \\
                           & \se{0.491} & \se{0.236}      & \se{0.363}       & \se{0.364}       \\
      \midrule
      Firm FE              & No  & Yes & Yes & Yes \\
      Product FE           & No  & No  & Yes & Yes \\
      City FE              & No  & No  & No  & Yes \\
      $N$                  & 28,957 & 28,626 & 28,277 & 28,277 \\
      Adj $R^2$            & 0.352  & 0.476  & 0.598  & 0.595  \\
      \bottomrule
    \end{tabular}
  \end{threeparttable}
\end{table}



% ------------------------------------------------------------
\subsection{交互项}
% ------------------------------------------------------------

\subsubsection*{回归方程}

\begin{equation}
  \ln p_{fjct} = \gamma_3 \ln Q_{fjt} + \boldsymbol{\lambda} \cdot \mathbf{z}_{jct}
    + \sum_{k} \mu_k \bigl(\ln\mathrm{Output}_{ft} \times z^k_{jct}\bigr)
    + \delta_f + \delta_j + \delta_c + \varepsilon_{fjct}
\end{equation}


\subsubsection*{变量说明}

\begin{table}[H]
  \centering\small
  \begin{tabular}{L{4.8cm} L{9.2cm}}
    \toprule
    变量 & 说明 \\
    \midrule
    $\ln\mathrm{Output} \times \ln n^B$       & 规模$\times$买方厚 \\[3pt]
    $\ln\mathrm{Output} \times \ln n^S$       & 规模$\times$卖方竞争 \\[3pt]
    $\ln\mathrm{Output} \times \ln\bar{p}^{B}$  & 规模$\times$买方侧市场均价 \\[3pt]
    $\ln\mathrm{Output} \times \ln\bar{p}^{S}$  & 规模$\times$卖方侧市场均价 \\[3pt]
    $\ln\mathrm{Output} \times \ln q^m$       & 规模$\times$市场规模 \\
    \bottomrule
  \end{tabular}
\end{table}

\subsubsection*{关键发现}

\begin{itemize}[nosep]
  \item 全部为全 FE 规格，$\ln\mathrm{Output}$ 仅通过交互项进入、主效应不单独估计。
  \item 单项交互揭示规模异质性：大企业在买方/卖方更多的市场中价格更低（$\mu_1=-0.019^{***}$，$\mu_2=-0.023^{***}$）；规模$\times$买方/卖方侧均价均为正（$+0.011^{***}$），即大企业对市场均价更敏感（此时对应主效应转负）。
  \item 全部交互同时加入后因多重共线性大多不显著，仅 Size$\times q^m$（$-0.011^{***}$）稳健；产量 $\ln Q$ 在全部列稳健显著为负。
\end{itemize}



\subsubsection*{回归结果}

\begin{table}[H]
  \centering
  \caption{T6：企业规模与市场条件交互项（全FE规格）}
  \begin{threeparttable}
    \begin{tabular}{lcccccc}
      \toprule
       & (1) Baseline & (2) +Size$\times n^B$ & (3) +Size$\times n^S$
         & (4) +Size$\times\bar{p}^{B}$ & (5) +Size$\times\bar{p}^{S}$ & (6) All \\
      \midrule
      $\ln Q$      & $-0.047\sym{***}$ & $-0.047\sym{***}$ & $-0.047\sym{***}$ & $-0.047\sym{***}$ & $-0.047\sym{***}$ & $-0.046\sym{***}$ \\
                   & \se{0.006}        & \se{0.006}        & \se{0.006}        & \se{0.006}        & \se{0.006}        & \se{0.006} \\[4pt]
      $\ln n^B$    & $0.133\sym{***}$ & $0.516\sym{***}$ & $0.138\sym{***}$ & $0.138\sym{***}$ & $0.138\sym{***}$ & $0.009$ \\
                   & \se{0.040}       & \se{0.101}       & \se{0.040}       & \se{0.040}       & \se{0.040}       & \se{0.148} \\[4pt]
      $\ln n^S$    & $-0.003$ & $-0.010$ & $0.457\sym{***}$ & $-0.005$ & $-0.006$ & $0.211$  \\
                   & \se{0.035} & \se{0.035} & \se{0.115}       & \se{0.034} & \se{0.035} & \se{0.156} \\[4pt]
      $\ln q^m$    & $-0.083\sym{***}$ & $-0.084\sym{***}$ & $-0.085\sym{***}$ & $-0.087\sym{***}$ & $-0.085\sym{***}$ & $0.134\sym{*}$ \\
                   & \se{0.017}        & \se{0.017}        & \se{0.017}        & \se{0.017}        & \se{0.017}        & \se{0.072} \\[4pt]
      $\ln\bar{p}^{B}$ & $-0.018$ & $-0.018$ & $-0.019$ & $-0.236\sym{***}$ & $-0.017$ & $0.083$ \\
                     & \se{0.017} & \se{0.017} & \se{0.017} & \se{0.073} & \se{0.017} & \se{0.087} \\[4pt]
      $\ln\bar{p}^{S}$ & $0.052\sym{***}$ & $0.052\sym{***}$ & $0.052\sym{***}$ & $0.052\sym{***}$ & $-0.161\sym{**}$ & $-0.075$ \\
                     & \se{0.011} & \se{0.011} & \se{0.011} & \se{0.011} & \se{0.064} & \se{0.074} \\[4pt]
      Size$\times n^B$        & \mis & $-0.019\sym{***}$ & \mis & \mis & \mis & $0.007$ \\
                              &     & \se{0.005}        &     &     &     & \se{0.007} \\[4pt]
      Size$\times n^S$        & \mis & \mis & $-0.023\sym{***}$ & \mis & \mis & $-0.011$ \\
                              &     &     & \se{0.006}        &     &     & \se{0.008} \\[4pt]
      Size$\times\bar{p}^{B}$ & \mis & \mis & \mis & $0.011\sym{***}$ & \mis & $-0.005$ \\
                              &     &     &     & \se{0.004}       &     & \se{0.004} \\[4pt]
      Size$\times\bar{p}^{S}$ & \mis & \mis & \mis & \mis & $0.011\sym{***}$ & $0.006$ \\
                              &     &     &     &     & \se{0.003}       & \se{0.004} \\[4pt]
      Size$\times q^m$        & \mis & \mis & \mis & \mis & \mis & $-0.011\sym{***}$ \\
                              &     &     &     &     &     & \se{0.004} \\[4pt]
      Cons       & $4.749\sym{***}$ & $4.702\sym{***}$ & $4.726\sym{***}$ & $4.786\sym{***}$ & $4.773\sym{***}$ & $4.744\sym{***}$ \\
                   & \se{0.298}       & \se{0.291}       & \se{0.290}       & \se{0.299}       & \se{0.299}       & \se{0.290}       \\
      \midrule
      Firm FE      & Yes & Yes & Yes & Yes & Yes & Yes \\
      Product FE   & Yes & Yes & Yes & Yes & Yes & Yes \\
      City FE      & Yes & Yes & Yes & Yes & Yes & Yes \\
      $N$          & 45,669 & 45,669 & 45,669 & 45,669 & 45,669 & 45,669 \\
      Adj $R^2$    & 0.605  & 0.605  & 0.605  & 0.605  & 0.605  & 0.606  \\
      \bottomrule
    \end{tabular}
  \end{threeparttable}
\end{table}



% ============================================================
\section{结果与模型预期的差距}
% ============================================================

本节对照模型预期，梳理实证结果中"符合预期"与"与预期有差距"的部分，
并指出最值得进一步处理的识别问题。

\subsection*{符合预期的部分}

\begin{itemize}[nosep]
  \item \textbf{数量折扣}：产量 $\ln Q$ 在所有规格下均稳健显著为负
    （OLS $-0.177^{***}$，全 FE $-0.047^{***}$），与"采购量越大、单价越低"的
    数量折扣预期一致。
  \item \textbf{规模$\times$市场厚度的交互为负}：大企业在买方更多
    （$\mu_1=-0.019^{***}$）、卖方更多（$\mu_2=-0.023^{***}$）以及市场总量更大
    （Size$\times q^m=-0.011^{***}$）的市场中支付更低单价，符合"规模赋予议价能力"
    的预期。
  \item \textbf{卖方侧均价的独立信息}：卖方侧 leave-one-out 均价
    $\ln\bar{p}^{S}$ 在全 FE 下仍显著为正（$0.052^{***}$），与"同侧市场价格水平
    传导到企业采购价"的预期一致。
\end{itemize}

\subsection*{与预期有差距的部分}

\begin{itemize}[nosep]
  \item \textbf{企业规模主效应方向相反（最核心的偏离）}：模型预期规模越大、议价能力
    越强、采购单价越低；但 OLS 中 $\ln\mathrm{Output}$ 显著为正
    （$+0.234^{***}$），方向与预期相反。值得注意的是，规模与市场厚度的交互项
    \emph{却是负的}，说明在同一市场内部，大企业确实把价格压低；为正的只是规模的
    \emph{主效应}（跨市场比较）。因此该正号更可能反映质量／产品定位的选择效应
    （大企业采购更高质量投入），而非议价能力——$\ln\mathrm{Output}$ 不应被解读为
    议价能力的代理。
  \item \textbf{产品相似度方向与预期相反}：理论预期相似度越高、企业对该产品越熟悉
    （技术邻近、信息不对称更小、甚至可自产），采购单价应更低；但 $S_{mj}$、$C_{mj}$
    在 OLS 强显著为正（$+1.491^{***}$、$+1.481^{***}$），加入 Product FE 后大幅衰减
    至勉强显著（$0.186^{*}$、$0.263^{*}$）。一方面，正相关约 87\% 在 Product FE 下被
    吸收，说明它主要来自跨产品的选择构成，而非"同一产品内、相似度更高则更便宜"；
    另一方面，残留的微弱正效应更可能反映定制化／质量溢价（与自身结构高度匹配的往往
    是更专用、更高规格的变体），与企业规模主效应为正属同一类"质量／定位"故事。
  \item \textbf{卖方数不显著}：模型预期卖方（供给）越多、竞争越激烈、单价越低，但
    $\ln n^S$ 在主回归中不显著，未呈现预期的负向竞争效应。
  \item \textbf{买方数为正}：$\ln n^B$ 全规格稳健为正（$0.115^{***}$），与"需求越多
    议价越弱"的朴素预期方向相反；若成立需重新解释为需求拥挤／市场景气效应。
  \item \textbf{市场总量符号不稳}：$\ln q^m$ 在 OLS 为正（$+0.077^{***}$），加入
    Product FE 后转为显著为负（$-0.100^{***}$），符号随规格翻转，说明其捕捉的成分
    在控制产品构成前后并不一致。
  \item \textbf{买方侧均价的塌缩}：$\ln\bar{p}^{B}$ 在 OLS/Firm FE 强正，但加入
    Product FE 后归零，这属于机械相关／产品构成被吸收，而非模型本身的失败。
\end{itemize}

\subsection*{小结与处理方向}

核心偏离在于\emph{企业规模主效应的方向}，而非交互项——交互项的负号说明规模在市场
内部仍带来议价优势。因此 $\ln\mathrm{Output}$ 的正主效应更可能刻画的是质量／产品
定位，而非议价能力。后续可考虑：(a) 引入质量代理变量（如单位价值、产品档次）以分离
质量效应；(b) 在文字中重新界定 $\ln\mathrm{Output}$ 的角色；(c) 在固定产品质量档次
后做稳健性检验，检验规模主效应是否仍为正。

\paragraph{数据口径提示。}
需特别说明的是，当前价格 $p_{fjct}$ 由发票金额除以数量得到，而\emph{数量的计量单位
尚未统一}（不同企业、不同商品可能以千克、吨、个、箱等不同单位申报）。单位不一致会使
同一 9 位商品下的单价不可直接比较，给上述系数（尤其是水平方向的符号与量级）带来噪声
干扰。后续应在统一计量单位、或在更细的"商品$\times$单位"层面重新计算单价后复核，以确认
本节的偏离是经济机制所致而非口径问题。

% ============================================================
\section{SQL 取数逻辑}
% ============================================================

本节说明从增值税数据库中提取分析所需数据的核心逻辑，包括五个CSV视图及主要Python处理步骤。

\subsection{五个核心CSV视图}

\begin{enumerate}[nosep, leftmargin=*]

  \item \textbf{样本企业采购汇总（firm\_buy）}：购方企业ID $\times$ 商品编码，含金额与数量字段；原始行数473,268行，清洗后保留425,713行，涉及2,731种商品。

  \item \textbf{样本企业销售汇总（firm\_sell）}：销方企业ID $\times$ 商品编码，含金额与数量字段；原始行数101,844行，清洗后保留89,584行，涉及2,542种商品。销售数量用于构造产量 $\ln Q_{fjt}$ 与卖方侧 LOO 的本企业扣除项。

  \item \textbf{企业地区对照（firm\_city）}：企业ID $\times$ 地级市，共3,410行。

  \item \textbf{全量城市采购汇总（city\_buy）}：购方地区 $\times$ 商品编码，含买方企业数、金额、数量字段；原始行数2,104,901行，清洗后保留1,393,329行。用于构造 $n^B_{jct}$、$q^m_{jct}$ 与 $\bar{p}^{B}_{jct}$。

  \item \textbf{全量城市销售汇总（city\_sell）}：销方地区 $\times$ 商品编码，含卖方企业数、金额、数量字段；原始行数1,540,065行，清洗后保留954,222行。用于构造 $n^S_{jct}$ 与 $\bar{p}^{S}_{jct}$。

\end{enumerate}

\subsection{核心Python处理步骤}

\begin{enumerate}[nosep, leftmargin=*]

  \item \textbf{商品码标准化}：将原始商品编码右补零至19位后取前9位，仅保留2,778个合法编码。

  \item \textbf{外包产品识别}：计算 firm\_buy $\cap$ firm\_sell 的交集，要求 $\min(\text{采购额},\,\text{销售额}) > 0$，识别为外包产品。

  \item \textbf{企业地区合并}：通过 firm\_city 为企业层面数据补充城市字段，用于后续城市固定效应及市场条件匹配。

  \item \textbf{市场条件构造}：基于全量口径（city\_buy 与 city\_sell）聚合，不限于样本企业，确保 $n^B$、$n^S$、$q^m$、$\bar{p}^{B}$、$\bar{p}^{S}$ 代表完整市场信息。两个 leave-one-out 均价在取比值前先从城市-产品总量中减去本企业自身的买方/卖方金额与数量（本企业为唯一买方/卖方时设为缺失）。

  \item \textbf{剔除中介企业}：依据企业特征数据中的标记（外包比例$>90\%$），剔除外包企业面板（2,108家）中的233家中介企业，最终回归面板包含\textbf{1,875家企业}。

\end{enumerate}

\subsection{样本汇总}

\begin{table}[H]
  \centering
  \caption{样本基本统计}
  \begin{tabular}{lc}
    \toprule
    指标 & 数值 \\
    \midrule
    原始抽样企业数                               & 3,410 \\
    采购端有9位产品码企业数                  & 3,376 \\
    外包面板企业数& 2,108 \\
    \quad（外包定义：同一产品既购又售）          & \\
    回归面板企业数    & 1,875 \\
    \quad（剔除233家中介企业后）                & \\
    平均外包产品种数/企业（均值/中位数）         & 25.0 / 3 \\
    外包产品总种数（跨所有企业）                 & 2,143 \\
    回归面板观测数（T3/T6，全FE）               & $\approx$45,669 \\
    涵盖城市数                                   & 262 \\
    \bottomrule
  \end{tabular}
\end{table}

% ============================================================
\section{桥接数据来源与匹配率}
% ============================================================

\subsection{背景与桥接路径}

增值税数据库中不含企业财务信息，无法直接获取资产总额（Capital）。为纳入 $\ln K_{ft}$ 变量，需通过汇算清缴数据桥接.

\subsection{匹配率统计}

\begin{table}[H]
  \centering
  \caption{企业层面桥接匹配率}
  \begin{tabular}{lc}
    \toprule
    指标 & 数值 \\
    \midrule
    回归面板企业总数（\texttt{reg\_panel.dta}）& 1,875 \\
    成功匹配企业数（有 $\ln K$）              & $\approx$1,168 \\
    企业层面匹配率                            & 62.3\% \\
    T1/T2/T4/T5有效观测数                    & $\approx$29,000 \\
    T3/T6有效观测数（全FE）                  & $\approx$45,669 \\
    因匹配缺失损失比例                        & $\approx$37\% \\
    \bottomrule
  \end{tabular}
\end{table}

\end{document}
